\documentclass[10pt,a4paper]{article}
\usepackage[margin=2cm]{geometry}
\usepackage[T1]{fontenc}
\usepackage[utf8]{inputenc}
\usepackage{lmodern}
\usepackage{microtype}
\usepackage{booktabs}
\usepackage{amsmath}
\usepackage{enumitem}
\usepackage{xcolor}
\usepackage[colorlinks=true,linkcolor=black,urlcolor=blue!50!black,citecolor=black]{hyperref}

\setlist{nosep,leftmargin=*}
\setlength{\parskip}{3pt}
\setlength{\parindent}{0pt}
\newcommand{\code}[1]{\texttt{\small #1}}

\title{\vspace{-1.2cm}\textbf{CS7IS3 Assignment 1: A Lucene Search Engine for the Cranfield Collection}}
\author{Name: \rule{4cm}{0.4pt} \qquad Student ID: \rule{3cm}{0.4pt}}
\date{\vspace{-0.6cm}}

\begin{document}
\maketitle

\section{Implementation}

The system is a Java~21/Maven project built on Apache Lucene~10.3.2
(\code{lucene-core}, \code{lucene-analysis-common}, \code{lucene-queryparser}).
A single script, \code{run.sh}, builds the jar, compiles NIST \code{trec\_eval}, indexes the
collection, runs every analyzer/similarity combination and prints a summary table.
\code{run.sh demo} starts an interactive search mode that prints the parsed query and the top~10 titles.

\begin{itemize}
  \item \textbf{Parsing} (\code{CranParser}). Reads the SMART-format files
  \code{cran.all.1400} and \code{cran.qry}, splitting records on \code{.I} and collecting the
  \code{.T} (title), \code{.A} (author), \code{.B} (bibliography) and \code{.W} (abstract)
  sections. The query ids in \code{cran.qry} are not contiguous (001, 002, 004, \dots, 365),
  while \code{cranqrel} numbers the queries 1--225 in file order. Queries are therefore
  renumbered by position; without this step most evaluation scores would be wrong.
  \item \textbf{Indexing} (\code{Indexer}). Each abstract becomes a Lucene \code{Document}
  with a stored, untokenised \code{id} (\code{StringField}) and \code{TextField}s for
  \code{title}, \code{author}, \code{biblio} and \code{content}. There is one index per
  analyzer, so analyzers can be compared side by side.
  \item \textbf{Searching} (\code{Searcher}). Each query is escaped and parsed by a
  \code{MultiFieldQueryParser} over \code{title} and \code{content} with OR semantics, using
  the same analyzer as at index time. The scoring model is set on the \code{IndexSearcher};
  the built-in similarities share the same length-norm encoding, so one index serves every
  model. The top 1000 documents per query are written in TREC run format
  (\code{qid Q0 docno rank score tag}).
  \item \textbf{Evaluation}. Cleverdon's grades run from 1 (complete answer) to 4 (minimum
  interest), but \code{trec\_eval} treats higher values as more relevant, so grades are mapped
  $g \mapsto 5-g$. The 225 entries graded $-1$ are not real judgements (the same document
  appears for several unrelated queries) and are dropped, leaving 1612 judgements. Runs are
  scored with MAP, P@5, P@10, Recall@100, Recall@1000 and nDCG@10.
\end{itemize}

\section{Choice of Analyzers}

Five analysis chains were compared (Table~\ref{tab:analyzers}). The chosen \emph{custom}
analyzer is:
\code{StandardTokenizer} $\rightarrow$ \code{EnglishPossessiveFilter} $\rightarrow$
\code{LowerCaseFilter} $\rightarrow$ \code{ASCIIFoldingFilter} $\rightarrow$
\code{StopFilter} $\rightarrow$ \code{PorterStemFilter}.

\begin{itemize}
  \item \textbf{Tokenisation.} \code{StandardTokenizer} follows Unicode word-break rules. It
  removes punctuation such as the frequent ``\code{ .}'' and splits hyphenated compounds such
  as \emph{boundary-layer} into their parts, so they match unhyphenated query terms.
  Whitespace splitting keeps punctuation attached to words and was the worst chain.
  \item \textbf{Stop words.} The Cranfield queries are full natural-language questions
  (\emph{``what similarity laws must be obeyed when constructing aeroelastic models\dots''}).
  Many query words are question words, auxiliaries or filler (\emph{what, can, does, obtained,
  available, so far}). They are missing from Lucene's 33-word default list, yet occur in most
  abstracts. Because each query is an OR of its terms, they add noise to the scores, so about
  50 such words were added to the Lucene English stop list.
  \item \textbf{Stemming.} Aerodynamics abstracts use many inflected forms
  (\emph{plate/plates, heated/heating, flows/flow}). Stemming gave the largest single gain.
  Porter's aggressive conflation beat the more conservative Krovetz stemmer (KStem) slightly.
\end{itemize}

\section{Scoring Models}

\begin{itemize}
  \item \textbf{Vector Space Model}: \code{ClassicSimilarity}, Lucene's TF-IDF cosine model
  with $tf=\sqrt{f}$, $idf = 1+\ln\frac{N}{df+1}$ and length norm $1/\sqrt{|d|}$.
  \item \textbf{BM25}: \code{BM25Similarity} with $k_1=1.2$, $b=0.75$. A grid over
  $k_1\in\{0.6,1.2,2.0\}$ and $b\in\{0.5,0.75,0.9\}$ changed MAP by less than 0.012, so the
  defaults were kept.
  \item \textbf{Additional models}: a Dirichlet-smoothed language model ($\mu=1000$), a
  Jelinek--Mercer language model ($\lambda=0.7$), and Divergence From Randomness (DFR,
  I(n)-B-H2).
\end{itemize}

\textbf{Field weighting.} A title boost of 2 was tried first and then tuned. The \code{.W}
abstract already repeats the title, so extra title weight double-counts it. With custom
analyzer + BM25, MAP was 0.305 (title ignored), \textbf{0.317 (boost 1)}, 0.301 (boost 2) and
0.281 (boost 4). Boost 1 is used in every reported run.

\section{Results}

\begin{table}[h]
\centering\small
\caption{Scoring models with the custom analyzer (225 queries, top 1000 results per query).}
\label{tab:sims}
\begin{tabular}{lcccccc}
\toprule
Similarity & MAP & P@5 & P@10 & R@100 & R@1000 & nDCG@10\\
\midrule
DFR (I(n)-B-H2)            & \textbf{0.3327} & \textbf{0.3502} & \textbf{0.2573} & \textbf{0.7669} & \textbf{0.9499} & \textbf{0.3974}\\
BM25                       & 0.3172 & 0.3316 & 0.2484 & 0.7482 & 0.9486 & 0.3854\\
LM Jelinek--Mercer         & 0.3051 & 0.3218 & 0.2400 & 0.7381 & 0.9496 & 0.3770\\
VSM (Classic TF-IDF)       & 0.2860 & 0.3022 & 0.2262 & 0.7249 & 0.9486 & 0.3554\\
LM Dirichlet               & 0.2792 & 0.2942 & 0.2147 & 0.7161 & 0.9497 & 0.3443\\
\bottomrule
\end{tabular}
\end{table}

\begin{table}[h]
\centering\small
\caption{Effect of the analyzer on MAP.}
\label{tab:analyzers}
\begin{tabular}{llccc}
\toprule
Analyzer & Pipeline & VSM & BM25 & DFR\\
\midrule
whitespace & whitespace split only                         & 0.2108 & 0.2535 & 0.2716\\
standard   & StandardTokenizer + lowercase                 & 0.2309 & 0.2848 & 0.3043\\
english    & Lucene \code{EnglishAnalyzer} (33 stop words, Porter) & 0.2838 & 0.3107 & 0.3297\\
kstem      & custom stop list + KStem                      & 0.2838 & 0.3135 & 0.3274\\
custom     & custom stop list + Porter                     & \textbf{0.2860} & \textbf{0.3172} & \textbf{0.3327}\\
\bottomrule
\end{tabular}
\end{table}

\textbf{Discussion.} The probabilistic models clearly beat the vector space model: BM25
improves MAP by 11\% relative to TF-IDF (0.317 vs.\ 0.286), and DFR does best on every
precision measure (MAP 0.333, P@5 0.350). BM25 saturates term frequency and normalises
for document length. This suits short abstracts, where a repeated term should not dominate
the score. DFR has no parameters to tune and is robust on small collections. The Dirichlet
language model is the weakest probabilistic model: with $\mu=1000$ and abstracts averaging
about 160 words, the collection statistics outweigh each document's own term statistics.
Analysis matters as much as the scoring model. Under BM25, going from whitespace tokenisation
to the custom chain raises MAP from 0.254 to 0.317 (+25\%). Most of this comes from correct
tokenisation and lowercasing (+0.031) and from stemming and stop words (+0.026); the extended
stop list adds a further +0.007. Recall@1000 is about 0.95 for every configuration, so the
models differ in how they \emph{rank} relevant documents, not in whether they find them.
BM25 is used for the demonstration because it is the standard Lucene baseline. DFR is
reported as the best-performing configuration.

\begin{thebibliography}{9}\small
\setlength{\itemsep}{0pt}
\bibitem{lucene} Apache Software Foundation. \emph{Apache Lucene 10.3 Documentation}. \url{https://lucene.apache.org/}
\bibitem{cleverdon} C.~Cleverdon. The Cranfield tests on index language devices. \emph{Aslib Proceedings}, 19(6):173--194, 1967.
\bibitem{salton} G.~Salton, A.~Wong, C.~S.~Yang. A vector space model for automatic indexing. \emph{Communications of the ACM}, 18(11):613--620, 1975.
\bibitem{bm25} S.~Robertson, H.~Zaragoza. The probabilistic relevance framework: BM25 and beyond. \emph{Foundations and Trends in IR}, 3(4):333--389, 2009.
\bibitem{dfr} G.~Amati, C.~J.~van~Rijsbergen. Probabilistic models of information retrieval based on measuring the divergence from randomness. \emph{ACM TOIS}, 20(4):357--389, 2002.
\bibitem{lm} C.~Zhai, J.~Lafferty. A study of smoothing methods for language models applied to information retrieval. \emph{ACM TOIS}, 22(2):179--214, 2004.
\bibitem{porter} M.~F.~Porter. An algorithm for suffix stripping. \emph{Program}, 14(3):130--137, 1980.
\bibitem{trec} NIST. \emph{trec\_eval}. \url{https://github.com/usnistgov/trec_eval}
\end{thebibliography}

\end{document}
